\begin{thebibliography}{99}
\bibitem{nice2022} National Institute for Health and Care Excellence. Fetal monitoring in labour. NICE guideline NG229. London, NICE, 2022, updated 25 March 2026. \url{https://www.nice.org.uk/guidance/ng229} [accessed 6 September 2026].
\bibitem{figo2015} Ayres-de-Campos D, Spong CY, Chandraharan E. FIGO consensus guidelines on intrapartum fetal monitoring: cardiotocography. Int J Gynaecol Obstet 2015;131:13--24. \href{https://doi.org/10.1016/j.ijgo.2015.06.020}{https://doi.org/10.1016/j.ijgo.2015.06.020}.
\bibitem{alfirevic2017} Alfirevic Z, Devane D, Gyte GML, Cuthbert A. Continuous cardiotocography (CTG) as a form of electronic fetal monitoring (EFM) for fetal assessment during labour. Cochrane Database Syst Rev 2017;2:CD006066. \href{https://doi.org/10.1002/14651858.CD006066.pub3}{https://doi.org/10.1002/14651858.CD006066.pub3}.
\bibitem{engelhart2023} Hernandez Engelhart C, Gundro Brurberg K, Aanstad KJ, et al. Reliability and agreement in intrapartum fetal heart rate monitoring interpretation: a systematic review. Acta Obstet Gynecol Scand 2023;102:970--85. \href{https://doi.org/10.1111/aogs.14591}{https://doi.org/10.1111/aogs.14591}.
\bibitem{infant2017} INFANT Collaborative Group. Computerised interpretation of fetal heart rate during labour (INFANT): a randomised controlled trial. Lancet 2017;389:1719--29. \href{https://doi.org/10.1016/S0140-6736(17)30568-8}{https://doi.org/10.1016/S0140-6736(17)30568-8}.
\bibitem{francis2024} Francis F, Luz S, Wu H, Stock SJ, Townsend R. Machine learning on cardiotocography data to classify fetal outcomes: a scoping review. Comput Biol Med 2024;172:108220. \href{https://doi.org/10.1016/j.compbiomed.2024.108220}{https://doi.org/10.1016/j.compbiomed.2024.108220}.
\bibitem{aeberhard2023} Aeberhard JL, Radan AP, Delgado-Gonzalo R, et al. Artificial intelligence and machine learning in cardiotocography: a scoping review. Eur J Obstet Gynecol Reprod Biol 2023;281:54--62. \href{https://doi.org/10.1016/j.ejogrb.2022.12.008}{https://doi.org/10.1016/j.ejogrb.2022.12.008}.
\bibitem{park2025} Park CE, Choi B, Park RW, et al. Automated interpretation of cardiotocography using deep learning in a nationwide multicenter study. Sci Rep 2025;15:19617. \href{https://doi.org/10.1038/s41598-025-02849-4}{https://doi.org/10.1038/s41598-025-02849-4}.
\bibitem{pegorie2024} Pegorie C, Liu B, Thilaganathan B, Bhide A. Antenatal noninvasive fetal electrocardiography: a literature review. Matern Fetal Med 2024;6:178--89. \href{https://doi.org/10.1097/FM9.0000000000000237}{https://doi.org/10.1097/FM9.0000000000000237}.
\bibitem{garcia2018} Garcia-Casado J, Ye-Lin Y, Prats-Boluda G, et al. Electrohysterography in the diagnosis of preterm birth: a review. Physiol Meas 2018;39:02TR01. \href{https://doi.org/10.1088/1361-6579/aaad56}{https://doi.org/10.1088/1361-6579/aaad56}.
\bibitem{patel2025} Patel S, Conway AE, Adjei T, et al. Is it possible to monitor fetal movements with a wearable device? A review of novel technologies. Eur J Obstet Gynecol Reprod Biol 2025;305:329--38. \href{https://doi.org/10.1016/j.ejogrb.2024.12.011}{https://doi.org/10.1016/j.ejogrb.2024.12.011}.
\bibitem{liu2024} Liu L, Pu Y, Fan J, et al. Wearable sensors, data processing, and artificial intelligence in pregnancy monitoring: a review. Sensors 2024;24:6426. \href{https://doi.org/10.3390/s24196426}{https://doi.org/10.3390/s24196426}.
\bibitem{levance2026} Le Vance J, Adeoye A, Man R, et al. Remote home cardiotocography: a systematic review and meta-analysis. PLOS Digit Health 2026;5:e0001184. \href{https://doi.org/10.1371/journal.pdig.0001184}{https://doi.org/10.1371/journal.pdig.0001184}.
\bibitem{stricker2025} Stricker K, Radan AP, Surbek D. Continuous remote home monitoring solutions for mother and fetus: a scoping review. Eur J Obstet Gynecol Reprod Biol 2025;305:170--7. \href{https://doi.org/10.1016/j.ejogrb.2024.12.018}{https://doi.org/10.1016/j.ejogrb.2024.12.018}.
\bibitem{lekadir2025} Lekadir K, Frangi AF, Porras AR, et al. FUTURE-AI: international consensus guideline for trustworthy and deployable artificial intelligence in healthcare. BMJ 2025;388:e081554. \href{https://doi.org/10.1136/bmj-2024-081554}{https://doi.org/10.1136/bmj-2024-081554}.
\bibitem{kilinc2025} Kılınçdemir Turgut Ü. Artificial intelligence and perinatology: a study on accelerated academic production---a bibliometric analysis. Front Med 2025;12:1505450. \href{https://doi.org/10.3389/fmed.2025.1505450}{https://doi.org/10.3389/fmed.2025.1505450}.
\bibitem{sibanda2024} Sibanda K, Ndayizigamiye P, Twinomurinzi H. Industry 4.0 technologies in maternal health care: bibliometric analysis and research agenda. JMIR Pediatr Parent 2024;7:e47848. \href{https://doi.org/10.2196/47848}{https://doi.org/10.2196/47848}.
\bibitem{xiao2024} Xiao J, Shen F, Lu W, Yu Z, Li S, Wu J. Identifying and visualizing global research trends and hotspots of artificial intelligence in medical ultrasound: a bibliometric analysis. Curr Med Imaging 2024;20:e15734056324388. \href{https://doi.org/10.2174/0115734056324388240919112351}{https://doi.org/10.2174/0115734056324388240919112351}.
\bibitem{xie2024} Xie W, Cai P, Hu Y, et al. AI-driven paradigm shift in computerized cardiotocography analysis: a systematic review and promising directions. Neurocomputing 2024;607:128446. \href{https://doi.org/10.1016/j.neucom.2024.128446}{https://doi.org/10.1016/j.neucom.2024.128446}.
\bibitem{malekifar2025} Malekifar M, Cao S, Phan K, Vo K, Zeng FG, Cao H. A survey of deep learning for non-invasive fetal electrocardiogram analysis. IEEE Access 2025;13:202553--75. \href{https://doi.org/10.1109/ACCESS.2025.3621737}{https://doi.org/10.1109/ACCESS.2025.3621737}.
\bibitem{huang2026} Huang ZZ, Zhang WT, Li Y, Hu HB. A review of deep learning methods for non-invasive fetal electrocardiography. Biomed Signal Process Control 2026;120:110124. \href{https://doi.org/10.1016/j.bspc.2026.110124}{https://doi.org/10.1016/j.bspc.2026.110124}.
\bibitem{donthu2021} Donthu N, Kumar S, Mukherjee D, Pandey N, Lim WM. How to conduct a bibliometric analysis: an overview and guidelines. J Bus Res 2021;133:285--96. \href{https://doi.org/10.1016/j.jbusres.2021.04.070}{https://doi.org/10.1016/j.jbusres.2021.04.070}.
\bibitem{page2021} Page MJ, McKenzie JE, Bossuyt PM, et al. The PRISMA 2020 statement: an updated guideline for reporting systematic reviews. BMJ 2021;372:n71. \href{https://doi.org/10.1136/bmj.n71}{https://doi.org/10.1136/bmj.n71}.
\bibitem{rethlefsen2021} Rethlefsen ML, Kirtley S, Waffenschmidt S, et al. PRISMA-S: an extension to the PRISMA statement for reporting literature searches in systematic reviews. Syst Rev 2021;10:39. \href{https://doi.org/10.1186/s13643-020-01542-z}{https://doi.org/10.1186/s13643-020-01542-z}.
\bibitem{aria2017} Aria M, Cuccurullo C. bibliometrix: an R-tool for comprehensive science mapping analysis. J Informetr 2017;11:959--75. \href{https://doi.org/10.1016/j.joi.2017.08.007}{https://doi.org/10.1016/j.joi.2017.08.007}.
\bibitem{bradford1934} Bradford SC. Sources of information on specific subjects. Engineering 1934;137:85--6.
\bibitem{lotka1926} Lotka AJ. The frequency distribution of scientific productivity. J Wash Acad Sci 1926;16:317--23. \url{https://www.jstor.org/stable/24529203}.
\bibitem{marx2014} Marx W, Bornmann L, Barth A, Leydesdorff L. Detecting the historical roots of research fields by reference publication year spectroscopy (RPYS). J Assoc Inf Sci Technol 2014;65:751--64. \href{https://doi.org/10.1002/asi.23089}{https://doi.org/10.1002/asi.23089}.
\bibitem{vaneck2010} van Eck NJ, Waltman L. Software survey: VOSviewer, a computer program for bibliometric mapping. Scientometrics 2010;84:523--38. \href{https://doi.org/10.1007/s11192-009-0146-3}{https://doi.org/10.1007/s11192-009-0146-3}.
\bibitem{small1973} Small H. Co-citation in the scientific literature: a new measure of the relationship between two documents. J Am Soc Inf Sci 1973;24:265--9. \href{https://doi.org/10.1002/asi.4630240406}{https://doi.org/10.1002/asi.4630240406}.
\bibitem{kessler1963} Kessler MM. Bibliographic coupling between scientific papers. Am Doc 1963;14:10--25. \href{https://doi.org/10.1002/asi.5090140103}{https://doi.org/10.1002/asi.5090140103}.
\bibitem{vaneck2009} van Eck NJ, Waltman L. How to normalize cooccurrence data? An analysis of some well-known similarity measures. J Am Soc Inf Sci Technol 2009;60:1635--51. \href{https://doi.org/10.1002/asi.21075}{https://doi.org/10.1002/asi.21075}.
\bibitem{pons2005} Pons P, Latapy M. Computing communities in large networks using random walks. In: Yolum P, Güngör T, Gürgen F, Özturan C, editors. Computer and Information Sciences, ISCIS 2005. Lecture Notes in Computer Science, vol. 3733. Berlin, Springer, 2005, p. 284--93. \href{https://doi.org/10.1007/11569596_31}{https://doi.org/10.1007/11569596\_31}.
\bibitem{callon1991} Callon M, Courtial JP, Laville F. Co-word analysis as a tool for describing the network of interactions between basic and technological research: the case of polymer chemistry. Scientometrics 1991;22:155--205. \href{https://doi.org/10.1007/BF02019280}{https://doi.org/10.1007/BF02019280}.
\bibitem{blondel2008} Blondel VD, Guillaume JL, Lambiotte R, Lefebvre E. Fast unfolding of communities in large networks. J Stat Mech 2008;2008:P10008. \href{https://doi.org/10.1088/1742-5468/2008/10/P10008}{https://doi.org/10.1088/1742-5468/2008/10/P10008}.
\bibitem{warrick2010} Warrick PA, Hamilton EF, Precup D, Kearney RE. Classification of normal and hypoxic fetuses from systems modeling of intrapartum cardiotocography. IEEE Trans Biomed Eng 2010;57:771--9. \href{https://doi.org/10.1109/TBME.2009.2035818}{https://doi.org/10.1109/TBME.2009.2035818}.
\bibitem{ferrario2009} Ferrario M, Signorini MG, Magenes G. Complexity analysis of the fetal heart rate variability: early identification of severe intrauterine growth-restricted fetuses. Med Biol Eng Comput 2009;47:911--9. \href{https://doi.org/10.1007/s11517-009-0502-8}{https://doi.org/10.1007/s11517-009-0502-8}.
\bibitem{fergus2016} Fergus P, Idowu I, Hussain A, Dobbins C. Advanced artificial neural network classification for detecting preterm births using EHG records. Neurocomputing 2016;188:42--9. \href{https://doi.org/10.1016/j.neucom.2015.01.107}{https://doi.org/10.1016/j.neucom.2015.01.107}.
\bibitem{spilka2017} Spilka J, Frecon J, Leonarduzzi R, Pustelnik N, Abry P, Doret M. Sparse support vector machine for intrapartum fetal heart rate classification. IEEE J Biomed Health Inform 2017;21:664--71. \href{https://doi.org/10.1109/JBHI.2016.2546312}{https://doi.org/10.1109/JBHI.2016.2546312}.
\bibitem{fergus2018} Fergus P, Selvaraj M, Chalmers C. Machine learning ensemble modelling to classify caesarean section and vaginal delivery types using cardiotocography traces. Comput Biol Med 2018;93:7--16. \href{https://doi.org/10.1016/j.compbiomed.2017.12.002}{https://doi.org/10.1016/j.compbiomed.2017.12.002}.
\bibitem{comert2018} Cömert Z, Kocamaz AF, Subha V. Prognostic model based on image-based time-frequency features and genetic algorithm for fetal hypoxia assessment. Comput Biol Med 2018;99:85--97. \href{https://doi.org/10.1016/j.compbiomed.2018.06.003}{https://doi.org/10.1016/j.compbiomed.2018.06.003}.
\bibitem{comert2019} Cömert Z, Kocamaz AF. Fetal hypoxia detection based on deep convolutional neural network with transfer learning approach. In: Software Engineering and Algorithms in Intelligent Systems. Adv Intell Syst Comput 2019;763:239--48. \href{https://doi.org/10.1007/978-3-319-91186-1_25}{https://doi.org/10.1007/978-3-319-91186-1\_25}.
\bibitem{ponsiglione2021} Ponsiglione AM, Cosentino C, Cesarelli G, Amato F, Romano M. A comprehensive review of techniques for processing and analyzing fetal heart rate signals. Sensors 2021;21:6136. \href{https://doi.org/10.3390/s21186136}{https://doi.org/10.3390/s21186136}.
\bibitem{georgoulas2017} Georgoulas G, Karvelis P, Spilka J, Chudáček V, Stylios CD, Lhotská L. Investigating pH based evaluation of fetal heart rate (FHR) recordings. Health Technol 2017;7:241--54. \href{https://doi.org/10.1007/s12553-017-0201-7}{https://doi.org/10.1007/s12553-017-0201-7}.
\bibitem{krupa2022} Krupa AJD, Dhanalakshmi S, Lai KW, Tan YQ, Wu X. An IoMT enabled deep learning framework for automatic detection of fetal QRS: a solution to remote prenatal care. J King Saud Univ Comput Inf Sci 2022;34:7200--11. \href{https://doi.org/10.1016/j.jksuci.2022.07.002}{https://doi.org/10.1016/j.jksuci.2022.07.002}.
\bibitem{chudacek2014} Chudáček V, Spilka J, Burša M, et al. Open access intrapartum CTG database. BMC Pregnancy Childbirth 2014;14:16. \href{https://doi.org/10.1186/1471-2393-14-16}{https://doi.org/10.1186/1471-2393-14-16}.
\bibitem{signorini2003} Signorini MG, Magenes G, Cerutti S, Arduini D. Linear and nonlinear parameters for the analysis of fetal heart rate signal from cardiotocographic recordings. IEEE Trans Biomed Eng 2003;50:365--74. \href{https://doi.org/10.1109/TBME.2003.808824}{https://doi.org/10.1109/TBME.2003.808824}.
\bibitem{goldberger2000} Goldberger AL, Amaral LA, Glass L, et al. PhysioBank, PhysioToolkit, and PhysioNet: components of a new research resource for complex physiologic signals. Circulation 2000;101:e215--20. \href{https://doi.org/10.1161/01.CIR.101.23.e215}{https://doi.org/10.1161/01.CIR.101.23.e215}.
\bibitem{felezorz2008} Fele-\v{Z}or\v{z} G, Kav\v{s}ek G, Novak-Antoli\v{c} \v{Z}, Jager F. A comparison of various linear and non-linear signal processing techniques to separate uterine EMG records of term and pre-term delivery groups. Med Biol Eng Comput 2008;46:911--22. \href{https://doi.org/10.1007/s11517-008-0350-y}{https://doi.org/10.1007/s11517-008-0350-y}.
\bibitem{chawla2002} Chawla NV, Bowyer KW, Hall LO, Kegelmeyer WP. SMOTE: synthetic minority over-sampling technique. J Artif Intell Res 2002;16:321--57. \href{https://doi.org/10.1613/jair.953}{https://doi.org/10.1613/jair.953}.
\bibitem{benmbarek2026} Ben M'Barek I, Ben M'Barek B, Jauvion G, et al. Randomised study of human machine collaboration for cardiotocography interpretation during labour. NPJ Digit Med 2026;9:365. \href{https://doi.org/10.1038/s41746-026-02556-y}{https://doi.org/10.1038/s41746-026-02556-y}.
\bibitem{cruzrivera2020} Cruz Rivera S, Liu X, Chan AW, et al. Guidelines for clinical trial protocols for interventions involving artificial intelligence: the SPIRIT-AI extension. BMJ 2020;370:m3210. \href{https://doi.org/10.1136/bmj.m3210}{https://doi.org/10.1136/bmj.m3210}.
\bibitem{liu2020consort} Liu X, Cruz Rivera S, Moher D, et al. Reporting guidelines for clinical trial reports for interventions involving artificial intelligence: the CONSORT-AI extension. BMJ 2020;370:m3164. \href{https://doi.org/10.1136/bmj.m3164}{https://doi.org/10.1136/bmj.m3164}.
\bibitem{tamber2020} Tamber KK, Hayes DJL, Carey SJ, Wijekoon JHB, Heazell AEP. A systematic scoping review to identify the design and assess the performance of devices for antenatal continuous fetal monitoring. PLoS One 2020;15:e0242983. \href{https://doi.org/10.1371/journal.pone.0242983}{https://doi.org/10.1371/journal.pone.0242983}.
\bibitem{mhajna2020} Mhajna M, Schwartz N, Levit-Rosen L, et al. Wireless, remote solution for home fetal and maternal heart rate monitoring. Am J Obstet Gynecol MFM 2020;2:100101. \href{https://doi.org/10.1016/j.ajogmf.2020.100101}{https://doi.org/10.1016/j.ajogmf.2020.100101}.
\bibitem{zhao2019} Zhao Z, Deng Y, Zhang Y, et al. DeepFHR: intelligent prediction of fetal acidemia using fetal heart rate signals based on convolutional neural network. BMC Med Inform Decis Mak 2019;19:286. \href{https://doi.org/10.1186/s12911-019-1007-5}{https://doi.org/10.1186/s12911-019-1007-5}.
\bibitem{fischer2023} Fischer AM, Rietveld AL, Teunissen PW, Bakker PCAM, Hoogendoorn M. End-to-end learning with interpretation on electrohysterography data to predict preterm birth. Comput Biol Med 2023;158:106846. \href{https://doi.org/10.1016/j.compbiomed.2023.106846}{https://doi.org/10.1016/j.compbiomed.2023.106846}.
\bibitem{huang2025} Huang ZZ, Zhang WT, Li Y, Cui J, Zhang YR. TCGAN: temporal convolutional generative adversarial network for fetal ECG extraction using single-channel abdominal ECG. IEEE J Biomed Health Inform 2025;29:3478--87. \href{https://doi.org/10.1109/JBHI.2024.3524085}{https://doi.org/10.1109/JBHI.2024.3524085}.
\bibitem{vandewiele2021} Vandewiele G, Dehaene I, Kovács G, et al. Overly optimistic prediction results on imbalanced data: a case study of flaws and benefits when applying over-sampling. Artif Intell Med 2021;111:101987. \href{https://doi.org/10.1016/j.artmed.2020.101987}{https://doi.org/10.1016/j.artmed.2020.101987}.
\bibitem{riley2024} Riley RD, Archer L, Snell KIE, et al. Evaluation of clinical prediction models (part 2): how to undertake an external validation study. BMJ 2024;384:e074820. \href{https://doi.org/10.1136/bmj-2023-074820}{https://doi.org/10.1136/bmj-2023-074820}.
\bibitem{collins2024} Collins GS, Moons KGM, Dhiman P, et al. TRIPOD+AI statement: updated guidance for reporting clinical prediction models that use regression or machine learning methods. BMJ 2024;385:e078378. \href{https://doi.org/10.1136/bmj-2023-078378}{https://doi.org/10.1136/bmj-2023-078378}.
\bibitem{moons2025} Moons KGM, Damen JAA, Kaul T, et al. PROBAST+AI: an updated quality, risk of bias, and applicability assessment tool for prediction models using regression or artificial intelligence methods. BMJ 2025;388:e082505. \href{https://doi.org/10.1136/bmj-2024-082505}{https://doi.org/10.1136/bmj-2024-082505}.
\bibitem{vasey2022} Vasey B, Nagendran M, Campbell B, et al. Reporting guideline for the early stage clinical evaluation of decision support systems driven by artificial intelligence: DECIDE-AI. BMJ 2022;377:e070904. \href{https://doi.org/10.1136/bmj-2022-070904}{https://doi.org/10.1136/bmj-2022-070904}.
\bibitem{rieke2020} Rieke N, Hancox J, Li W, et al. The future of digital health with federated learning. NPJ Digit Med 2020;3:119. \href{https://doi.org/10.1038/s41746-020-00323-1}{https://doi.org/10.1038/s41746-020-00323-1}.
\end{thebibliography}
